\chapter{The Liouville Measure and Hamiltonian Flow}

The transition from the macroscopic description (thermodynamics) to the microscopic description (statistical mechanics) is framed as a move from the ``surface'' (the contact manifold) to the ``volume'' (the symplectic phase space). This module utilises the tools of measure theory to explain how the collective behaviour of microstates gives rise to the macroscopic Legendrian manifolds discussed in Module~I.

\section{The Symplectic Phase Space of Classical Mechanics}

In classical mechanics, the phase space $\Gamma$ of a system with $N$ degrees of freedom is a $2N$-dimensional smooth manifold equipped with a closed, non-degenerate 2-form $\SympForm$:

\begin{definition}[Symplectic manifold]
A \textbf{symplectic manifold} is a pair $(\Gamma^{2N}, \SympForm)$ where:
\begin{enumerate}
\item $\SympForm$ is \textbf{closed}: $\dd\SympForm = 0$.
\item $\SympForm$ is \textbf{non-degenerate}: $\SympForm^N \neq 0$ (i.e.\ the $N$-th exterior power is a volume form).
\end{enumerate}
\end{definition}

In canonical coordinates $(q^1, \ldots, q^N, p_1, \ldots, p_N)$, the symplectic form takes the standard Darboux form:
\begin{equation}
\SympForm = \sum_{i=1}^{N} \dd q^i \wedge \dd p_i.
\end{equation}

For a gas of $\mathcal{N}$ particles in three-dimensional space, $N = 3\mathcal{N}$, and the phase space is $6\mathcal{N}$-dimensional.

\section{The Liouville Volume Form}

The non-degeneracy of $\SympForm$ provides a canonical volume form on $\Gamma$:

\begin{definition}[Liouville measure]
The \textbf{Liouville volume form} is the $2N$-form:
\begin{equation}
\boxed{\mu_L = \frac{1}{N!}\,\SympForm^N = \frac{1}{N!}\,\underbrace{\SympForm \wedge \cdots \wedge \SympForm}_{N\text{ times}}.}
\end{equation}
In canonical coordinates:
\begin{equation}
\mu_L = \dd q^1 \wedge \dd p_1 \wedge \dd q^2 \wedge \dd p_2 \wedge \cdots \wedge \dd q^N \wedge \dd p_N.
\end{equation}
\end{definition}

This is the unique (up to normalisation) volume form that is determined by the symplectic structure alone --- no metric or additional structure is needed. The Liouville measure $\mu_L$ defines the natural ``counting measure'' on the phase space: the volume of a region $\Omega \subset \Gamma$ is:
\begin{equation}
\text{Vol}(\Omega) = \int_\Omega \mu_L.
\end{equation}

\section{Hamiltonian Flow and the Equations of Motion}

\begin{definition}[Hamiltonian vector field]
Given a smooth function $H: \Gamma \to \R$ (the Hamiltonian), the \textbf{Hamiltonian vector field} $X_H$ is defined by:
\begin{equation}
\iota_{X_H}\,\SympForm = -\dd H,
\end{equation}
where $\iota$ denotes the interior product (contraction).
\end{definition}

In canonical coordinates, this yields Hamilton's equations:
\begin{equation}
\dot{q}^i = \pdv{H}{p_i}, \qquad \dot{p}_i = -\pdv{H}{q^i}.
\end{equation}

The Hamiltonian flow $\phi_t: \Gamma \to \Gamma$ is the one-parameter family of diffeomorphisms generated by $X_H$: the map $\phi_t(q_0, p_0)$ gives the state of the system at time $t$ if it started at $(q_0, p_0)$ at time $0$.

\section{Liouville's Theorem}

The central result of this chapter is that Hamiltonian flow preserves the Liouville volume:

\begin{theorem}[Liouville's theorem]
The Hamiltonian flow $\phi_t$ preserves the Liouville volume form:
\begin{equation}
\boxed{\Lie_{X_H}\,\mu_L = 0,}
\end{equation}
where $\Lie_{X_H}$ is the Lie derivative along the Hamiltonian vector field.
\end{theorem}

\begin{proof}
By Cartan's magic formula:
\begin{equation}
\Lie_{X_H}\,\SympForm = \iota_{X_H}\,\dd\SympForm + \dd(\iota_{X_H}\,\SympForm).
\end{equation}
The first term vanishes because $\dd\SympForm = 0$ (closedness of $\SympForm$). The second term is:
\begin{equation}
\dd(\iota_{X_H}\,\SympForm) = \dd(-\dd H) = -\dd^2 H = 0,
\end{equation}
by the nilpotency of the exterior derivative. Therefore $\Lie_{X_H}\,\SympForm = 0$: the Hamiltonian flow preserves the symplectic form.

Since $\mu_L = \frac{1}{N!}\SympForm^N$, and the Lie derivative obeys the Leibniz rule on exterior products:
\begin{equation}
\Lie_{X_H}\,\mu_L = \frac{1}{N!}\sum_{k=1}^{N} \SympForm^{k-1} \wedge (\Lie_{X_H}\,\SympForm) \wedge \SympForm^{N-k} = 0.
\end{equation}
\end{proof}

\subsection{Physical interpretation}

Liouville's theorem states that the phase-space volume of any region is preserved under Hamiltonian evolution:
\begin{equation}
\text{Vol}(\phi_t(\Omega)) = \text{Vol}(\Omega) \qquad \forall\, t.
\end{equation}

This is the geometric statement that Hamiltonian mechanics is \emph{incompressible} in phase space. The ``fluid'' of representative points neither converges nor diverges --- it moves like an incompressible fluid.

\begin{remark}
Liouville's theorem is the reason we can define stationary probability distributions on phase space. Without the symplectic structure ensuring this volume invariance, the concept of a \emph{stationary ensemble} would be mathematically ill-defined: any probability density would be ``washed away'' by the flow.
\end{remark}

\section{The Liouville Equation}

Let $\rho(q, p, t)$ be a probability density on phase space. The requirement that probability is conserved under Hamiltonian flow leads to the \textbf{Liouville equation}:
\begin{equation}
\boxed{\pdv{\rho}{t} + \{H, \rho\} = 0,}
\end{equation}
where $\{H, \rho\}$ is the Poisson bracket:
\begin{equation}
\{H, \rho\} = \sum_{i=1}^{N}\left(\pdv{H}{q^i}\pdv{\rho}{p_i} - \pdv{H}{p_i}\pdv{\rho}{q^i}\right).
\end{equation}

Equivalently, using the Lie derivative:
\begin{equation}
\pdv{\rho}{t} + \Lie_{X_H}\rho = 0.
\end{equation}

A \textbf{stationary distribution} satisfies $\partial\rho/\partial t = 0$, which requires $\{H, \rho\} = 0$: the density must be constant along the Hamiltonian flow. Any function $\rho = f(H)$ satisfies this automatically.

\section{The Energy Hypersurface and the Microcanonical Measure}

For a system with a fixed energy $E$, the dynamics are confined to the \textbf{energy hypersurface}:
\begin{equation}
\Sigma_E = \{(q, p) \in \Gamma : H(q, p) = E\}.
\end{equation}

This is a $(2N-1)$-dimensional submanifold of $\Gamma$. The natural measure on $\Sigma_E$ is induced from the Liouville measure by the ``coarea formula'':
\begin{equation}
\dd\mu_{\Sigma_E} = \frac{\mu_L}{\|\dd H\|}\bigg|_{\Sigma_E}.
\end{equation}

This is the \textbf{microcanonical measure}: the uniform measure on the energy surface, weighted by the inverse gradient of the Hamiltonian. The density of states is:
\begin{equation}
\Omega(E) = \int_{\Sigma_E} \dd\mu_{\Sigma_E},
\end{equation}
and the microcanonical entropy is $S(E) = k_B \ln\Omega(E)$.

\section{Summary}

\begin{table}[H]
\centering
\renewcommand{\arraystretch}{1.6}
\begin{tabular}{ll}
\toprule
\textbf{Structure} & \textbf{Role in Statistical Mechanics} \\
\midrule
Symplectic form $\SympForm$ & Defines the geometry of phase space \\
Liouville measure $\mu_L = \frac{1}{N!}\SympForm^N$ & Canonical volume form (``counting measure'') \\
Hamiltonian flow $\phi_t$ & Time evolution of microstates \\
Liouville's theorem $\Lie_{X_H}\mu_L = 0$ & Volume preservation $\Rightarrow$ stationary ensembles exist \\
Energy surface $\Sigma_E$ & Domain of the microcanonical ensemble \\
\bottomrule
\end{tabular}
\end{table}

The symplectic structure of classical mechanics provides the Liouville measure, and Liouville's theorem guarantees that this measure is preserved under time evolution. This is the mathematical foundation upon which all statistical ensembles are built.
